% Original vector figure for Version 3.9.0.
\begin{figure}[htbp]
\centering
\begin{tikzpicture}[font=\small,>=Stealth,line width=.65pt]
\begin{scope}[xshift=-3cm]
\draw[->,gray] (-1.3,0)--(1.3,0);\draw[->,gray] (0,-1)--(0,1);
\draw[blue,thick] (0,0) ellipse(1 and .65);
\node[below] at (0,-1) {ellipse};
\end{scope}
\begin{scope}[xshift=.1cm]
\draw[->,gray] (-1.3,0)--(1.3,0);\draw[->,gray] (0,-1)--(0,1);
\draw[blue,thick] plot[domain=-.8:.8,samples=35] ({sqrt(.35+\x*\x)},\x);
\draw[blue,thick] plot[domain=-.8:.8,samples=35] ({-sqrt(.35+\x*\x)},\x);
\node[below] at (0,-1) {hyperbola};
\end{scope}
\begin{scope}[xshift=3.2cm]
\draw[->,gray] (-1.2,0)--(1.2,0);\draw[->,gray] (0,-.2)--(0,1.1);
\draw[blue,thick] plot[domain=-1:1,samples=35] (\x,{\x*\x});
\node[below] at (0,-1) {parabola};
\end{scope}
\end{tikzpicture}
\caption{Three regular conics. Their equations become local graphs wherever the corresponding partial derivative is nonzero; a vertical tangent requires solving for the other variable.}\label{fig:conics-quadrics}
\end{figure}
